\section{GIỚI THIỆU}

\subsection{Giới thiệu về Công ty Cổ phần Dịch vụ Công nghệ Tin học HPT}

Công ty Cổ phần Dịch vụ Công nghệ Tin học HPT (HPT) là một trong những doanh nghiệp công nghệ thông tin hàng đầu tại Việt Nam với hơn 30 năm hình thành và phát triển. HPT hoạt động trong nhiều lĩnh vực như tư vấn và triển khai giải pháp công nghệ thông tin, tích hợp hệ thống, phát triển phần mềm, cung cấp dịch vụ chuyển đổi số, điện toán đám mây, an toàn thông tin, quản trị hạ tầng CNTT, đào tạo và cung ứng nguồn nhân lực công nghệ thông tin.

Trong suốt quá trình phát triển, HPT đã xây dựng được đội ngũ chuyên gia và kỹ sư giàu kinh nghiệm, không ngừng nghiên cứu và tiếp cận các công nghệ hiện đại như điện toán đám mây (Cloud Computing), trí tuệ nhân tạo (Artificial Intelligence), dữ liệu lớn (Big Data), Internet vạn vật (IoT) và các giải pháp chuyển đổi số cho doanh nghiệp. Với năng lực chuyên môn cao cùng kinh nghiệm triển khai nhiều dự án quy mô lớn, HPT đã trở thành đối tác tin cậy của nhiều tổ chức, doanh nghiệp và cơ quan trọng yếu tại Việt Nam.

Bên cạnh năng lực công nghệ, HPT còn chú trọng xây dựng môi trường làm việc chuyên nghiệp, đề cao giá trị văn hóa doanh nghiệp với phương châm ``Nhân bản -- Hài hòa'', tạo điều kiện cho nhân viên phát triển chuyên môn và đóng góp vào sự phát triển bền vững của doanh nghiệp. Hệ thống quy trình quản lý dự án, quản lý chất lượng và phát triển sản phẩm được xây dựng theo các chuẩn quốc tế nhằm đảm bảo chất lượng dịch vụ và nâng cao hiệu quả triển khai dự án.

\subsection{Giới thiệu nội dung thực tập}

Trong thời gian thực tập tại HPT, sinh viên được tham gia tìm hiểu môi trường làm việc thực tế của doanh nghiệp, quy trình quản lý dự án theo chuẩn CMMI, mô hình phát triển Agile Scrum cũng như các hoạt động quản lý chất lượng, quản lý rủi ro và phối hợp với khách hàng trong quá trình triển khai dự án.

Bên cạnh việc tìm hiểu các quy trình quản lý dự án, sinh viên còn được tiếp cận với dự án thực tế trong lĩnh vực Internet of Things (IoT). Dự án hướng đến xây dựng hệ thống giám sát, quản lý và bảo trì máy lọc nước từ xa theo thời gian thực. Thông qua dự án này, sinh viên có cơ hội tìm hiểu kiến trúc hệ thống IoT, quy trình thu thập và xử lý dữ liệu từ thiết bị nhúng, cơ chế truyền thông giữa thiết bị và máy chủ, cũng như các công nghệ được sử dụng trong quá trình phát triển sản phẩm tại doanh nghiệp.

Đợt thực tập là cơ hội để vận dụng kiến thức đã học vào môi trường thực tế, đồng thời nâng cao kỹ năng chuyên môn, kỹ năng làm việc nhóm và hiểu rõ hơn về quy trình phát triển và triển khai các dự án công nghệ thông tin trong doanh nghiệp.
\section{NỘI DUNG THỰC TẬP}

\subsection{Thông tin đề tài thực tập}

\textbf{Tên đề tài:} Phát triển hệ thống IoT hỗ trợ giám sát, quản lý và bảo trì máy lọc nước theo thời gian thực.

\textbf{Ngôn ngữ lập trình sử dụng:}
\begin{itemize}
    \item C/C++
    \item C\#
    \item JavaScript
    \item SQL
\end{itemize}

\textbf{Framework và công nghệ sử dụng:}
\begin{itemize}
    \item Arduino Framework
    \item MQTTnet
    \item ASP.NET Core MVC
    \item MQTT Broker
    \item Hệ quản trị cơ sở dữ liệu SQL
\end{itemize}

\subsection{Mục tiêu thực tập}

Mục tiêu của đợt thực tập là tìm hiểu quy trình phát triển hệ thống IoT trong môi trường doanh nghiệp, nghiên cứu kiến trúc hệ thống máy lọc nước thông minh và tham gia xây dựng các thành phần của hệ thống từ thiết bị phần cứng đến nền tảng quản lý tập trung. Bên cạnh đó, sinh viên được tiếp cận với các công nghệ Web Backend, giao thức truyền thông MQTT và quy trình phát triển phần mềm thực tế.

\subsection{Nội dung công việc thực hiện}

Trong suốt thời gian thực tập, các công việc được thực hiện theo từng giai đoạn như sau:

\begin{table}[htbp]
    \centering
    \begin{tabularx}{\textwidth}{|>{\centering\arraybackslash}p{2cm}|X|}
        \hline
        \textbf{Tuần} & \textbf{Nội dung công việc} \\
        \hline
        1 & Tìm hiểu về công ty, môi trường làm việc, quy trình phát triển sản phẩm và định hướng thực tập. \\
        \hline
        2 & Nghiên cứu hệ thống hiện tại của dự án, tìm hiểu kiến trúc giải pháp, quy trình nghiệp vụ và mô hình IoT đang được áp dụng. \\
        \hline
        3 & Làm quen với hệ thống Web Backend, khảo sát kiến trúc dữ liệu và luồng xử lý dữ liệu trong hệ thống. \\
        \hline
        4 & Khảo sát phần cứng máy lọc nước, nghiên cứu sơ đồ điều khiển, các cảm biến được sử dụng và đánh giá khả năng tích hợp IoT vào thiết bị. \\
        \hline
        5 & Tìm hiểu giao thức MQTT, thực hiện kết nối thiết bị với MQTT Broker và xây dựng cơ chế truyền dữ liệu từ thiết bị lên hệ thống quản lý. \\
        \hline
        6 & Tham gia phát triển các chức năng thu thập, xử lý và lưu trữ dữ liệu từ máy lọc nước; xây dựng giao diện Dashboard hỗ trợ giám sát thiết bị theo thời gian thực. \\
        \hline
        7 & Kiểm thử hệ thống, đánh giá kết quả hoạt động, tối ưu và hoàn thiện các chức năng được giao. \\
        \hline
        8 & Tổng hợp kết quả thực hiện, hoàn thiện tài liệu kỹ thuật và báo cáo thực tập. \\
        \hline
    \end{tabularx}
\end{table}